% ---
% Experimentos complementarios (X1--X8) y estudios de caracterización del simulador (B1--B3).
% Sección metodológica previa al reporte de resultados.
% Entradas que lo respaldan: H-018, H-019, H-035, H-036, H-041, H-044, H-045, H-046, H-047, H-048, H-051, H-052, H-053, H-054, D-035, D-038.
% ---

\subsection{Experimentos complementarios}
\label{sec:experimentos_complementarios}

Con el propósito de evaluar la solidez de las conclusiones obtenidas en el protocolo principal frente a hipótesis
alternativas o variantes metodológicas que un lector u evaluador razonable plantearía, se diseñó una batería de ocho
experimentos complementarios (denominados X1 a X8). Estos experimentos no alteran el diseño experimental básico ni
los puntos de decisión congelados para el pipeline principal ($k = 12$, $reps = 1$, Arquitectura~B con vector de
parámetros $\boldsymbol\theta$ fijo con semilla 42 y particiones independientes por paciente), sino que someten a
prueba la sensibilidad de los hallazgos ante variaciones controladas en la selección de características, la dimensionalidad
del registro cuántico, la arquitectura de clasificación, la optimización variacional y el protocolo de evaluación.
La Tabla~\ref{tab:resumen_experimentos_complementarios} resume las características de diseño de cada experimento, y
los párrafos siguientes detallan su formulación metodológica. Los resultados cuantitativos derivados de estas pruebas se
reportan en el Capítulo~\ref{sec:resultados}.

% fuente: H-035, H-036, H-041, H-044, H-045, H-046, H-047, H-048, H-053, H-054, D-035, D-038
\begin{table}[htbp]
\centering
\small
\caption{Resumen de diseño de los experimentos complementarios X1 a X8.}
\label{tab:resumen_experimentos_complementarios}
\begin{adjustbox}{max width=\linewidth}
\begin{tabular}{p{0.06\linewidth} p{0.22\linewidth} p{0.28\linewidth} p{0.22\linewidth} p{0.18\linewidth} p{0.14\linewidth}}
\toprule
\textbf{Exp.} & \textbf{Nombre / Archivo} & \textbf{Pregunta metodológica} & \textbf{Variación vs. Fijo} & \textbf{Control} & \textbf{Uso de prueba} \\
\midrule
X1 & \texttt{X1\_Selection\_}\allowbreak\texttt{Information\_Loss} & ¿Cuánta información se pierde al reducir de 67 a 12 variables? & Varía subconjunto de variables; fijos modelos clásicos. & 67 variables y PCA con 12 componentes. & No (solo entrenamiento). \\
X2 & \texttt{X2\_Kernel\_}\allowbreak\texttt{Class\_Structure} & ¿Conservan los kernels cuánticos similitud intra-clase? & Varía escala $c$; fijos kernels de C4 y C5. & Kernel RBF clásico y permutación. & No (submuestra de entrenamiento). \\
X3 & \texttt{X3\_Fewer\_Qubits\_}\allowbreak\texttt{k8} & ¿Persiste la conclusión con un registro de $k = 8$ qubits? & Varía $k \in \{8, 12\}$; fijo protocolo de métricas. & Condiciones clásicas con 8 variables. & No en separabilidad; sensibilidad en M7. \\
X4 & \texttt{X4\_QSVM\_and\_}\allowbreak\texttt{67\_Features} & ¿Se pierde información en la codificación o en la lectura? & Varía clasificador (QSVM vs MLP) y variables ($k=12$ vs 67). & SVM clásico RBF y MLP de M7. & Sí (una vez por configuración). \\
X5 & \texttt{X5\_Pipeline\_}\allowbreak\texttt{Audit} & ¿Existe algún fallo de implementación o fuga de datos? & Varía implementación de componentes; fijos datos. & Implementaciones alternativas y código Qiskit. & Sí (solo verificación de partición). \\
X6 & \texttt{X6\_Searching\_}\allowbreak\texttt{for\_Advantage} & ¿Supera algún kernel cuántico al clásico en etiquetas reales? & Búsqueda anidada sobre familias cuánticas y clásicas. & Mejores modelos clásicos y etiquetas de Huang. & Sí (una vez por mejor modelo). \\
X7 & \texttt{X7\_Theta\_}\allowbreak\texttt{Draws} & ¿Introduce sesgo el sorteo único de $\boldsymbol\theta$? & Varía semilla de $\boldsymbol\theta$ (10 sorteos); fijo MLP. & Variabilidad de MLP y condiciones clásicas. & No (solo validación cruzada). \\
X8 & \texttt{X8\_Trained\_}\allowbreak\texttt{Theta} & ¿Mejora la representación al entrenar $\boldsymbol\theta$? & Varía optimización de $\boldsymbol\theta$ (entrenado vs fijo); fijo MLP. & Modelo con $\boldsymbol\theta$ fijo de M7 y C1 crudo. & Sí (una vez por subconjunto y mapa). \\
\bottomrule
\end{tabular}
\end{adjustbox}
\end{table}

% fuente: H-035
\textbf{Experimento X1: Pérdida de información en la selección de características.} Desarrollado en
\texttt{Code/}\allowbreak\texttt{experiments/}\allowbreak\texttt{X1\_Selection\_Information\_Loss.ipynb}, este análisis responde a la
interrogante de cuánta información diagnóstica y varianza global se descarta al reducir el conjunto radiómico de 67 a 12
variables en M4. En su diseño se compara la capacidad predictiva y la varianza retenida evaluando distintos subconjuntos
de características: las 67 originales, las 12 seleccionadas por el protocolo estándar, 12 variables reajustadas dentro de
cada pliegue de validación cruzada, las 12 mejores según la prueba F univariante sin filtro de redundancia, y conjuntos de
12 variables elegidas al azar. Como control se emplean modelos clásicos de regresión logística y bosques aleatorios ajustados
exclusivamente dentro de los pliegues de validación cruzada sobre el conjunto de entrenamiento, junto con una descomposición
PCA de 12 componentes como cota lineal de retención de varianza. Este experimento no utiliza el conjunto de prueba oficial.
Los resultados se analizan en la Sección~\ref{sec:datos_caracteristicas_seleccion}.

% fuente: H-036
\textbf{Experimento X2: Estructura de clase en matrices de kernel.} Implementado en
\texttt{Code/}\allowbreak\texttt{experiments/}\allowbreak\texttt{X2\_Kernel\_Class\_Structure.ipynb}, evalúa si dos lesiones de la misma
clase diagnóstica (benigna--benigna o maligna--maligna) exhiben mayor similitud de Hilbert que dos lesiones de clases opuestas,
y en qué medida dicha diferenciación depende del factor de contracción angular $c$. El diseño analiza la distribución de
similitudes por pares de lesiones sobre la submuestra estratificada de 200 casos de entrenamiento, calculando el área
bajo la curva ROC de pares frente a un barrido del parámetro de escala angular. Como control de referencia se utiliza el kernel
gaussiano RBF clásico con ancho de banda calibrado por la heurística de la mediana, contrastando la significancia estadística
frente a una prueba de permutación de 1\,000 iteraciones con etiquetas barajadas. El experimento prescinde del conjunto de prueba.
Sus resultados se presentan en la Sección~\ref{sec:separabilidad_representaciones}.

% fuente: H-041, D-035
\textbf{Experimento X3: Evaluación con registro reducido ($k = 8$ qubits).} Desarrollado en
\texttt{Code/}\allowbreak\texttt{experiments/}\allowbreak\texttt{X3\_Fewer\_Qubits\_k8.ipynb}, examina si la ausencia de ventaja cuántica
en separabilidad observada en M6 es atribuible a la concentración exponencial del espacio de Hilbert a $k = 12$ qubits. El diseño
mantiene idénticos los protocolos de evaluación de KTA, dimensión efectiva, diferencia geométrica regularizada $g$ y razón de Fisher,
variando la dimensión del registro cuántico a $k = 8$ qubits (lo que reduce los acoplamientos $ZZ$ de 66 a 28 bloques y los parámetros
del circuito variacional de 24 a 16). Se toman de forma determinista las 8 primeras características del ordenamiento de M4. Como control,
se evalúan las condiciones clásicas C1, C2 y C3 sobre las mismas 8 características de entrada. En la fase de separabilidad este
experimento no accede al conjunto de prueba oficial, mientras que su traslación a clasificación se evalúa en el análisis de sensibilidad
del Módulo~7. Los resultados se detallan en la Sección~\ref{sec:separabilidad_representaciones}.

% fuente: H-044, H-045, D-038
\textbf{Experimento X4: Máquinas de soporte vectorial cuánticas y codificación de 67 variables.} Ejecutado en
\texttt{Code/}\allowbreak\texttt{experiments/}\allowbreak\texttt{X4\_QSVM\_and\_67\_Features.ipynb}, busca determinar si la pérdida de
desempeño observada en M7 reside intrínsecamente en la codificación cuántica en el espacio de Hilbert o en la etapa de lectura
mediante expectativas locales $\langle Z_i \rangle$ bajo la Arquitectura~B. Para ello, se implementa una máquina de vectores de
soporte con kernel cuántico (QSVM) evaluada directamente sobre las matrices completas de fidelidad de C4 y C5. Asimismo, se evalúa si
la disponibilidad de las 67 características radiómicas completas permite formular kernels cuánticos con ventaja frente al límite
clásico, analizando codificaciones en amplitud sobre 7 qubits y codificaciones angulares en producto tensorial sobre 67 qubits. Se
contrasta directamente contra modelos SVM clásicos lineales y RBF, así como contra el clasificador MLP entrenado sobre las expectativas
de M7. Los modelos se optimizan mediante validación cruzada y se evalúan una única vez sobre el conjunto de prueba oficial.
Los resultados se reportan en la Sección~\ref{sec:qsvm_ventaja}.

% fuente: H-046
\textbf{Experimento X5: Auditoría independiente del pipeline experimental.} Documentado en
\texttt{Code/}\allowbreak\texttt{experiments/}\allowbreak\texttt{X5\_Pipeline\_Audit.ipynb}, responde a la necesidad metodológica de
descartar que la ausencia de ventaja cuántica provenga de anomalías de software, fugas de datos (*data leakage*) entre entrenamiento
y prueba, asimetrías de preprocesamiento o discrepancias numéricas en el cómputo de fidelidades. La auditoría consta de 25 pruebas
programadas independientes que reconstruyen particiones y diagnósticos directamente desde las fuentes maestras de metadatos,
contrastan la simulación de kernels mediante circuitos físicos frente a la formulación por vectores de estado, y verifican la
equivalencia de las implementaciones clasificatorias frente a bibliotecas estándar de la comunidad. El conjunto de prueba se
utiliza exclusivamente para verificar la integridad y aislamiento de las particiones, sin entrenamiento ni ajuste de parámetros.
Su análisis se discute en la Sección~\ref{sec:qsvm_ventaja}.

% fuente: H-047, H-048
\textbf{Experimento X6: Búsqueda exhaustiva de ventaja cuántica y contraste con etiquetas sintéticas.} Desarrollado en
\texttt{Code/}\allowbreak\texttt{experiments/}\allowbreak\texttt{X6\_Searching\_for\_Advantage.ipynb}, plantea una búsqueda anidada y
simétrica de configuraciones óptimas en ambos paradigmas. En el lado cuántico se exploran mapas $ZZ$ con entrelazamiento lineal y
completo, repeticiones $reps \in \{1, 2\}$, factores de escala angular y kernels proyectados; en el lado clásico se evalúan familias
gaussianas, lineales y polinomiales. Para discernir si la ausencia de ventaja obedece al protocolo experimental o a la relación
geométrica intrínseca de las etiquetas patológicas reales, se incorpora una prueba de contraste con etiquetas sintéticas generadas
según la metodología de Huang et al., alineadas por construcción con la métrica del kernel cuántico. La selección de modelos se efectúa
mediante validación cruzada anidada y la evaluación en el conjunto de prueba se restringe a una sola ejecución por contendiente global.
Los resultados se exponen en la Sección~\ref{sec:qsvm_ventaja}.

% fuente: H-053
\textbf{Experimento X7: Sensibilidad al vector de parámetros variacionales $\boldsymbol\theta$.} Implementado en
\texttt{Code/}\allowbreak\texttt{experiments/}\allowbreak\texttt{X7\_Theta\_Draws.ipynb}, evalúa si la fijación de $\boldsymbol\theta$
mediante un único sorteo aleatorio (semilla 42) en la Arquitectura~B introdujo un sesgo desfavorable para las representaciones
cuánticas C4 y C5. El diseño analiza diez sorteos pseudoaleatorios independientes del vector $\boldsymbol\theta$ (semillas 42 a 51),
manteniendo fijos la semilla del clasificador MLP, la arquitectura de red y los pliegues de validación cruzada de M7. Como control
se compara la variación inducida por el parámetro $\boldsymbol\theta$ frente a la variabilidad por semilla de inicialización del
MLP clásico y frente a los desempeños obtenidos por C1 y C3. Este análisis opera de forma estricta sobre el conjunto de entrenamiento
en validación cruzada sin tocar el conjunto de prueba. Sus resultados se reportan en la Sección~\ref{sec:rendimiento_clasificacion}.

% fuente: H-054
\textbf{Experimento X8: Optimización variacional conjunta de $\boldsymbol\theta$.} Realizado en
\texttt{Code/}\allowbreak\texttt{experiments/}\allowbreak\texttt{X8\_Trained\_Theta.ipynb}, examina si el entrenamiento supervisado
de los parámetros $\boldsymbol\theta$ en forma simultánea con los pesos del clasificador MLP permite cerrar la brecha de rendimiento
observada frente a las condiciones clásicas. En este esquema, el vector de expectativas alimenta una capa diferencial que actualiza
$\boldsymbol\theta$ mediante retropropagación sobre el vector de estado simulado en PyTorch con una tasa de aprendizaje dedicada. Se
analiza la sensibilidad ante la inicialización de $\boldsymbol\theta$ y se compara el modelo resultante frente a la condición fija
de M7, frente a la condición clásica C1 y frente al clasificador QSVM de X4. El conjunto de prueba oficial se consulta exactamente
una vez por modelo optimizado. Los resultados se analizan en las Secciones~\ref{sec:rendimiento_clasificacion} y~\ref{sec:restricciones_simulador}.

% fuente: H-018, H-019, H-051, H-052
\textbf{Estudios de caracterización del simulador cuántico (B1--B3).} De forma complementaria a los experimentos sobre datos radiómicos,
se llevaron a cabo tres estudios sistemáticos en \texttt{Code/benchmark/} para cuantificar las restricciones operacionales de la simulación
clásica de circuitos cuánticos en el contexto del proyecto. El estudio B1 (\texttt{B1\_Quantum\_Benchmark.ipynb}) midió la complejidad
estructural de los circuitos (profundidad, número de compuertas CX y conteo total de compuertas en función de $k$ y $reps$) y proyectó
el coste computacional del entrenamiento conjunto bajo la Arquitectura~A frente al esquema de precomputación de la Arquitectura~B.
No obstante, la sección de sensibilidad al número de disparos (*shots*) de B1 quedó formalmente invalidada debido a que la referencia
analítica empleada incorporaba ruido gaussiano no nulo inyectado por omisión de precisión cero en el estimador de alto nivel.
El estudio B2 (\texttt{B2\_Design\_Verifications.ipynb}) integró las comprobaciones empíricas que respaldaron las decisiones de diseño
metodológico: número de qubits y profundidad, arquitectura del \textit{embedding}, condiciones experimentales, cálculo de la \textit{geometric difference} y validación cruzada. Finalmente, el estudio B3 (\texttt{B3\_Sampling\_and\_Concentration.ipynb}) sustituyó y rectificó la caracterización
de disparos de B1 mediante un esquema de muestreo real con el estimador por muestreo de cadenas de bits de Aer, verificando el escalamiento
estadístico de error teórico $1/\sqrt{n}$ sin suelo de dispersión espuria hasta $2^{20}$ disparos, además de cuantificar la concentración
asintótica del vector de expectativas en función de la dimensionalidad $k$ y la profundidad del circuito. Sus hallazgos se integran en
las Secciones~\ref{sec:restricciones_simulador} y~\ref{sec:lecciones_implementacion}.
